% =====================================================
% CHAPITRE 5 : MODÉLISATION (CRISP-DM Phase 4)
% =====================================================
\chapter{Modélisation}

% ---------------------------------------------------
\section*{Introduction}
\addcontentsline{toc}{section}{Introduction}
% ---------------------------------------------------

Une fois les données préparées, vient la phase de construction des modèles. Plutôt que de
choisir d'emblée l'approche la plus sophistiquée, j'ai opté pour une progression logique :
commencer par un modèle simple servant de référence, complexifier si les performances le
justifient, puis compléter avec une perspective non supervisée. Cette stratégie permet
d'évaluer la valeur ajoutée réelle de chaque niveau de complexité et d'éviter un
sur-ingénierie inutile. Ce chapitre décrit les trois architectures développées et les
décisions qui ont guidé leur conception, à travers les notebooks
\texttt{03\_baseline\_models.ipynb} et \texttt{04\_autoencoder.ipynb}.

% ---------------------------------------------------
\section{Régression logistique}
% ---------------------------------------------------

\subsection{Principe et rôle dans le projet}

La régression logistique est la \textbf{baseline} --- le modèle de référence que tout
algorithme plus complexe doit surpasser pour justifier sa complexité. Elle produit une
probabilité d'appartenance à la classe fraude via la fonction sigmoïde :

\[
P(\text{fraude} \mid \mathbf{x}) = \sigma(\mathbf{w}^\top \mathbf{x} + b)
= \frac{1}{1 + e^{-(\mathbf{w}^\top \mathbf{x} + b)}}
\]

Sa faible capacité expressive est ici un avantage : si la régression logistique détecte
déjà une grande partie des fraudes, cela signifie que le signal est linéairement séparable,
ce qui simplifie l'ensemble de la chaîne. Dans le cas contraire, cela justifie le passage
à des modèles non linéaires.

\subsection{Hyperparamètres et configurations évaluées}

\begin{table}[H]
\centering
\caption{Hyperparamètres de la régression logistique}
\label{tab:lr_params}
\begin{tabular}{lll}
\toprule
\textbf{Paramètre} & \textbf{Valeur} & \textbf{Justification} \\
\midrule
\texttt{C}             & 0,1      & Régularisation forte (peu de fraudes) \\
\texttt{solver}        & lbfgs   & Adapté aux datasets de taille moyenne \\
\texttt{class\_weight} & balanced & Correction automatique du déséquilibre \\
\texttt{max\_iter}     & 1000    & Assure la convergence \\
\bottomrule
\end{tabular}
\end{table}

Deux configurations sont comparées :
\begin{itemize}
  \item \textbf{LR\_balanced} : entraînée sur \texttt{X\_train} avec
    \texttt{class\_weight='balanced'} ;
  \item \textbf{LR\_smote} : entraînée sur \texttt{X\_smote} (données équilibrées par
    SMOTE) sans pondération explicite.
\end{itemize}

\subsection{Réglage du seuil de décision}

À seuil par défaut ($0{,}5$), \texttt{LR\_balanced} maximise fortement le rappel mais au
prix d'une précision très faible. Une recherche de seuil sur l'ensemble de validation est
donc réalisée à partir des probabilités produites par le modèle. Les seuils optimaux retenus
sont \textbf{0,9986} pour \texttt{LR\_balanced} et \textbf{0,9621} pour \texttt{LR\_smote}.
Après ce recalibrage, \texttt{LR\_smote} devient la meilleure variante logistique avec un
rappel de \textbf{0,6923}, une précision de \textbf{0,7297} et un F1-score de
\textbf{0,7105} sur l'ensemble de test.

% ---------------------------------------------------
\section{Forêt aléatoire}
% ---------------------------------------------------

\subsection{Principe et avantages}

La forêt aléatoire constitue la \textbf{référence supervisée robuste} de ce projet. Son
principe repose sur la diversité : entraîner un grand nombre d'arbres de décision sur des
sous-échantillons distincts des données (\textit{bootstrap}), puis agréger leurs
prédictions. Cette diversité rend l'ensemble résistant au sur-apprentissage
(\textit{overfitting}) et capable de capturer des interactions non linéaires entre variables
que la régression logistique ne peut pas modéliser.

La prédiction finale est la \textbf{moyenne des probabilités} émises par chaque arbre
individuel, ce qui produit des scores bien calibrés.

\subsection{Hyperparamètres}

\begin{table}[H]
\centering
\caption{Hyperparamètres de la forêt aléatoire}
\label{tab:rf_params}
\begin{tabular}{lll}
\toprule
\textbf{Paramètre} & \textbf{Valeur} & \textbf{Justification} \\
\midrule
\texttt{n\_estimators}      & 300 & Stabilité des estimations \\
\texttt{max\_depth}         & 10  & Limite l'overfitting sur peu de fraudes \\
\texttt{min\_samples\_leaf} & 5   & Évite les feuilles trop petites \\
\texttt{class\_weight}      & balanced & Correction du déséquilibre \\
\bottomrule
\end{tabular}
\end{table}

Deux variantes sont comparées :
\begin{itemize}
  \item \textbf{RF\_balanced} : entraînée sur \texttt{X\_train} avec
    \texttt{class\_weight='balanced'} ;
  \item \textbf{RF\_smote} : entraînée sur \texttt{X\_smote}, sans pondération explicite.
\end{itemize}

\subsection{Réglage du seuil et importance des variables}

Les seuils optimaux retenus après recherche sur validation sont \textbf{0,6235} pour
\texttt{RF\_balanced} et \textbf{0,6291} pour \texttt{RF\_smote}. La variante
\texttt{RF\_smote} obtient les meilleures performances supervisées globales sur le test :
rappel de \textbf{0,7949}, précision de \textbf{0,8158} et F1-score de \textbf{0,8052}.

La forêt aléatoire permet également de calculer l'importance de chaque variable via la
réduction moyenne d'impureté (\textit{Mean Decrease in Impurity}). Sur
\texttt{RF\_smote}, les cinq variables les plus influentes sont :

\begin{table}[H]
\centering
\caption{Importance des variables (RF\_smote)}
\label{tab:feature_importance_rf}
\begin{tabular}{lr}
\toprule
\textbf{Variable} & \textbf{Importance} \\
\midrule
\texttt{balance\_diff\_orig}       & $\sim$0,385 \\
\texttt{hour}                      & $\sim$0,118 \\
\texttt{dest\_zero\_balance}       & $\sim$0,115 \\
\texttt{type\_TRANSFER}            & $\sim$0,089 \\
\texttt{log\_amount}               & $\sim$0,081 \\
\bottomrule
\end{tabular}
\end{table}

Cette hiérarchie confirme que le comportement du compte émetteur, la temporalité et la
nature du virement sont les signaux les plus discriminants --- en cohérence avec l'EDA.

% ---------------------------------------------------
\section{Autoencoder}
% ---------------------------------------------------

\subsection{Principe de détection par reconstruction}

L'autoencoder repose sur une idée intuitive : \textbf{apprendre à compresser puis à
reconstruire fidèlement les données normales}. Si une transaction légitime peut être
reproduite depuis un espace latent de petite dimension, c'est qu'elle s'inscrit dans des
patterns réguliers. Une transaction frauduleuse --- jamais vue à l'entraînement --- sera
en revanche mal reconstruite, trahissant son caractère inhabituel par une erreur de
reconstruction élevée.

\begin{definitionbox}
\textbf{Principe de l'autoencoder pour la détection d'anomalies :}
\begin{itemize}
  \item \textbf{Entraînement :} uniquement sur les transactions légitimes → le modèle
    apprend à reconstruire fidèlement les patterns normaux.
  \item \textbf{Inférence :} sur une transaction frauduleuse (pattern jamais vu),
    l'erreur de reconstruction est élevée car le modèle ne sait pas comment la reproduire.
  \item \textbf{Décision :} si l'erreur de reconstruction (MSE) dépasse un seuil
    $\theta$, la transaction est classée comme anomalie.
\end{itemize}
\end{definitionbox}

\subsection{Architecture du réseau}

L'architecture retenue est un autoencoder symétrique à trois couches dans l'encodeur et
trois dans le décodeur, avec un espace latent (\textit{bottleneck}) de dimension 4 :

\begin{table}[H]
\centering
\caption{Architecture de l'autoencoder ($14 \to 10 \to 7 \to [4] \to 7 \to 10 \to 14$)}
\label{tab:ae_architecture}
\begin{tabular}{llll}
\toprule
\textbf{Partie} & \textbf{Couche} & \textbf{Dimension} & \textbf{Activation} \\
\midrule
Entrée   & Input      & 14 & --- \\
Encodeur & Dense 1    & 10 & ReLU + BatchNorm + Dropout(0,2) \\
         & Dense 2    & 7  & ReLU + BatchNorm + Dropout(0,2) \\
         & Bottleneck & 4  & ReLU \\
Décodeur & Dense 3    & 7  & ReLU + BatchNorm + Dropout(0,2) \\
         & Dense 4    & 10 & ReLU + BatchNorm + Dropout(0,2) \\
Sortie   & Dense 5    & 14 & Linéaire \\
\bottomrule
\end{tabular}
\end{table}

Les techniques de régularisation employées sont :
\begin{itemize}
  \item \textbf{BatchNorm} (normalisation par lots) : stabilise l'apprentissage ;
  \item \textbf{Dropout} (taux 0,2) : réduit le sur-apprentissage ;
  \item \textbf{Régularisation L2} ($\lambda = 10^{-5}$) sur les poids denses.
\end{itemize}

\subsection{Entraînement}

\begin{table}[H]
\centering
\caption{Paramètres d'entraînement de l'autoencoder}
\label{tab:ae_training}
\begin{tabular}{ll}
\toprule
\textbf{Paramètre} & \textbf{Valeur} \\
\midrule
Données d'entraînement  & 139~818 transactions légitimes (0 fraude) \\
Validation interne      & 10\,\% du train normal \\
Fonction de perte       & MSE (Mean Squared Error) \\
Optimiseur              & Adam ($\alpha = 10^{-3}$) \\
Batch size              & 256 \\
Epochs max              & 100 (EarlyStopping, patience = 10) \\
ReduceLROnPlateau       & facteur 0,5, patience 5, min\_lr $10^{-6}$ \\
\bottomrule
\end{tabular}
\end{table}

L'entraînement s'arrête après \textbf{35 époques effectives} grâce à \textit{EarlyStopping}.
Le meilleur \texttt{val\_loss} observé est de \textbf{0,145024} et le temps d'entraînement
est de \textbf{45,2 secondes}. Le modèle comporte \textbf{664 paramètres} entraînables,
ce qui reste très léger au regard de la taille du corpus.

\subsection{Détermination du seuil de détection}

Une première borne est estimée par le \textbf{95\textsuperscript{e} percentile} de
l'erreur de reconstruction sur les transactions normales d'entraînement
($\theta_{p95} = 0{,}403750$). Ce seuil est utile comme repère initial, mais trop permissif
pour une décision finale.

Le seuil optimal est ensuite trouvé sur l'ensemble de validation (\texttt{X\_val}) en
balayant 200 valeurs entre le 50\textsuperscript{e} et le 99,9\textsuperscript{e}
percentile des erreurs de reconstruction. Pour chaque valeur, le F1-score est calculé.
Le seuil maximisant ce score est retenu : \textbf{$\theta = 1{,}687408$}.

\subsection{Visualisation de l'espace latent}

La projection PCA du bottleneck (dimension 4) en 2D permet de vérifier qualitativement
que les fraudes de validation tendent à se positionner \textbf{en périphérie des zones de
forte densité normale} --- confirmant la pertinence du mécanisme de reconstruction comme
score d'anomalie.

% ---------------------------------------------------
\section{Récapitulatif des configurations}
% ---------------------------------------------------

\begin{table}[H]
\centering
\caption{Récapitulatif des configurations de modèles construites}
\label{tab:models_recap}
\begin{tabular}{p{2.8cm}p{3cm}p{3cm}p{3.5cm}}
\toprule
\textbf{Configuration} & \textbf{Type} & \textbf{Données train}
  & \textbf{Gestion déséquilibre} \\
\midrule
LR\_balanced  & Supervisé linéaire   & \texttt{X\_train}        & \texttt{class\_weight} \\
LR\_smote     & Supervisé linéaire   & \texttt{X\_smote}        & SMOTE \\
RF\_balanced  & Supervisé ensembliste & \texttt{X\_train}       & \texttt{class\_weight} \\
RF\_smote     & Supervisé ensembliste & \texttt{X\_smote}       & SMOTE \\
Autoencoder   & Non supervisé        & \texttt{X\_train\_normal} & Données normales seules \\
\bottomrule
\end{tabular}
\end{table}

% ---------------------------------------------------
\section*{Conclusion}
\addcontentsline{toc}{section}{Conclusion}
% ---------------------------------------------------

Ce chapitre a décrit la conception et l'implémentation des cinq configurations de modèles
étudiées. Les variantes supervisées montrent que le choix du seuil de décision influence
fortement le compromis précision/rappel, et que \texttt{RF\_smote} s'impose comme la
meilleure configuration supervisée. L'autoencoder complète ce dispositif par une logique
non supervisée fondée sur la reconstruction, particulièrement utile lorsque les labels
sont rares. Les premiers résultats sur l'ensemble de validation suggèrent que la forêt
aléatoire dépasse les objectifs fixés (recall $> 0{,}70$, F1 $> 0{,}70$). L'évaluation
comparative détaillée sur l'ensemble de test, ainsi que les mécanismes d'explicabilité,
sont présentés dans le chapitre suivant.
